% Prose rendering of PrimeTensor/Tensor/Contract.lean.
% Fragment: no \documentclass, no preamble, no document environment.

\section{Multiplicative contraction}
\label{Tensor:Contract:sec}

\leanfile{PrimeTensor/Tensor/Contract.lean}

\subsection*{What this file establishes}

The multiplicative trace of a square matrix: the product of its diagonal
entries.  Where the classical trace sums the diagonal, this multiplies it,
which is the right move for a carrier whose primitive operation is
multiplication and which offers no addition.

The file is three declarations, and its interest is that it needs no empty
case.  A product over an index set must return something when that set is
empty, and the only uniform answer is the unit of the operation --- so a
contraction defined over $\name{Fin}\;n$ would have to demand a unit from the
coefficient type and single out $n = 0$.  Here the dimension is a depth and
therefore at least one, the fold that computes the product needs no identity
element, and the coefficient type is asked only for a multiplication.

\subsection{The definition}

\begin{definition}[Contraction, \lean{Tensor.contract₂}]
\label{Tensor:Contract:def:contract}
For a type $A$ with a multiplication and
$m : \name{Matrix}\;A\;\mathit{dim}$,
\[
  \name{contract}_2(m) \defeq
    \name{fold}_{\,\cdot}\bigl(\mathit{dim},\;
      \lambda i.\; m.\name{component}\;(i,i)\bigr) \;:\; A.
\]
\end{definition}

\begin{leancode}
def contract₂ {A : Type} [Mul A] {dim : Depth} (m : Matrix A dim) : A :=
  Axis.fold (· * ·) dim (fun i => m.component (i, i))
\end{leancode}

Two things from earlier modules are being used without ceremony.  The
diagonal selector $\lambda i.\; m.\name{component}\;(i,i)$ is well typed
because at rank two the index type \emph{is} a product of two axis types,
definitionally, as established in
\texttt{PrimeTensor/Tensor/Basic.lean}.  And
$\name{Axis.fold}$, from \texttt{PrimeTensor/Depth.lean}, folds a function on
$\name{Axis}\;\mathit{dim}$ with a binary operation and no seed value.

\begin{proposition}[What the contraction is]\label{Tensor:Contract:prop:expansion}
Let $n = \size{\mathit{dim}}$ and let $i_1, \dots, i_n$ enumerate
$\name{Axis}\;\mathit{dim}$ in order.  Writing
$m_{jk} = m.\name{component}\;(j,k)$,
\[
  \name{contract}_2(m)
    = m_{i_1 i_1} \cdot
      \bigl( m_{i_2 i_2} \cdot
        \bigl( \cdots \cdot m_{i_n i_n} \bigr) \bigr),
\]
the right-nested product of the diagonal entries in axis order, with no unit
at the innermost position.
\end{proposition}

\begin{proof}
Immediate from Definition~\ref{Tensor:Contract:def:contract} and the expansion
of $\name{Axis.fold}$ proved in \texttt{PrimeTensor/Depth.lean}, applied to
the diagonal selector.
\end{proof}

\begin{designnote}[Only a magma, and the bracketing is real]
\label{Tensor:Contract:note:magma}
The instance argument is \lean{[Mul A]} alone: no associativity, no
commutativity, no unit.  Two consequences deserve stating, because for the
classical trace neither arises.

First, the bracketing in
Proposition~\ref{Tensor:Contract:prop:expansion} is part of the value, not an
artefact of how it is written.  Without associativity, a differently bracketed
product of the same diagonal entries need not be the same element of $A$.

Second, the \emph{order} of the diagonal entries is likewise part of the
value, fixed by the enumeration of $\name{Axis}\;\mathit{dim}$.  Without
commutativity there is no permutation invariance to appeal to.  For the
intended carriers --- strictly positive reals under multiplication --- both
laws do hold, so nothing is lost there; the definition simply does not assume
them, and so applies to carriers where they fail.
\end{designnote}

\subsection{The two computation rules}

\begin{lemma}[Base case, \lean{contract₂\_dim\_one}]
\label{Tensor:Contract:lem:one}
For $m : \name{Matrix}\;A\;\ctor{one}$,
\[
  \name{contract}_2(m) = m.\name{component}\;(\ctor{first}, \ctor{first}).
\]
\end{lemma}

\begin{lemma}[Successor case, \lean{contract₂\_dim\_succ}]
\label{Tensor:Contract:lem:succ}
For $m : \name{Matrix}\;A\;(\ctor{succ}\,d)$,
\[
  \name{contract}_2(m) =
    m.\name{component}\;(\ctor{first}, \ctor{first}) \cdot
      \name{fold}_{\,\cdot}\bigl(d,\;
        \lambda i.\; m.\name{component}\;
          (\ctor{next}\,i, \ctor{next}\,i)\bigr).
\]
\end{lemma}

\begin{leancode}
@[simp] theorem contract₂_dim_one {A : Type} [Mul A]
    (m : Matrix A .one) :
    contract₂ m = m.component (.first, .first) := rfl

@[simp] theorem contract₂_dim_succ {A : Type} [Mul A]
    {d : Depth} (m : Matrix A (.succ d)) :
    contract₂ m =
      m.component (.first, .first) *
        Axis.fold (· * ·) d (fun i => m.component (.next i, .next i)) := rfl
\end{leancode}

\begin{proof}
Both are $\ctor{rfl}$.  Unfolding
Definition~\ref{Tensor:Contract:def:contract} turns each side into an
application of $\name{Axis.fold}$ at a dimension in constructor form, where
the corresponding defining equation of $\name{fold}$ --- $\name{fold\_one}$
and $\name{fold\_succ}$ in \texttt{PrimeTensor/Depth.lean} --- applies
definitionally.  As there, the statements are content-free as propositions and
the \lean{@[simp]} tags are the point: they let the simplifier evaluate a
contraction whose dimension is a concrete numeral, which unfolding
$\name{contract}_2$ alone would not do.  Together they are the complete case
analysis on the dimension, so no third rule is possible.
\end{proof}

\begin{remark}[The successor rule is not recursive]
\label{Tensor:Contract:rem:not-recursive}
Lemma~\ref{Tensor:Contract:lem:succ} exposes a raw
$\name{Axis.fold}$ on its right-hand side rather than a smaller contraction,
and it has to: $m$ has type $\name{Matrix}\;A\;(\ctor{succ}\,d)$, whose
component function is defined on $\name{Axis}\;(\ctor{succ}\,d)$, and there is
no matrix of dimension $d$ sitting inside it to contract.

One can of course build the restriction along $\ctor{next}$,
\[
  m' \defeq \bigl\langle\, \lambda \mathit{ij}.\;
    m.\name{component}\;(\ctor{next}\,\mathit{ij}_1,
      \ctor{next}\,\mathit{ij}_2) \,\bigr\rangle
  \;:\; \name{Matrix}\;A\;d,
\]
and then $\name{contract}_2(m')$ is definitionally the fold appearing in the
lemma, so the rule could equally have been phrased recursively.  The file does
not introduce $m'$.  The practical consequence is that an induction on the
dimension about $\name{contract}_2$ is carried out at the level of
$\name{Axis.fold}$, generalising over the selector function, rather than by
recursion on contractions.  ($m'$ is not a declaration of the file; it is
written here only to say what the successor rule amounts to.)
\end{remark}

\begin{remark}[Scope]\label{Tensor:Contract:rem:scope}
Only rank two is treated, as the subscript in the name records: there is no
contraction of a general tensor, no contraction of a chosen pair of indices,
and no interaction with the outer product of
\texttt{PrimeTensor/Tensor/Basic.lean} --- in particular nothing here relates
$\name{contract}_2(\name{outer}(u,v))$ to the coefficients of $u$ and $v$.
Nor is anything proved about $\name{contract}_2$ beyond its two computation
rules.
\end{remark}

\subsection*{Summary of the correspondence}

\begin{center}
\small
\begin{tabular}{@{}lll@{}}
\hline
Lean declaration & Kind & Here \\
\hline
\lean{Tensor.contract₂}            & definition & Def.~\ref{Tensor:Contract:def:contract} \\
\lean{Tensor.contract₂\_dim\_one}  & simp thm   & Lem.~\ref{Tensor:Contract:lem:one} \\
\lean{Tensor.contract₂\_dim\_succ} & simp thm   & Lem.~\ref{Tensor:Contract:lem:succ} \\
\hline
\end{tabular}
\end{center}

\noindent
Every declaration of \texttt{PrimeTensor/Tensor/Contract.lean} appears above.
The file contains no \lean{sorry}, no axiom, and no unproved named
proposition; its two theorems are proved by $\ctor{rfl}$.
Proposition~\ref{Tensor:Contract:prop:expansion} and
Remark~\ref{Tensor:Contract:rem:not-recursive} are stated for the reader and
are not declarations of the file.
